\documentclass[11pt]{article}
\usepackage[margin=1in]{geometry}
\usepackage{amsmath,amssymb,amsthm}
\usepackage[hidelinks]{hyperref}
\newtheorem{theorem}{Theorem}
\newtheorem{proposition}[theorem]{Proposition}
\newtheorem{lemma}[theorem]{Lemma}
\newtheorem{corollary}[theorem]{Corollary}
\theoremstyle{remark}
\newtheorem{remark}[theorem]{Remark}
\newcommand{\F}{\mathbb F}
\newcommand{\E}{\mathbb E}
\newcommand{\PP}{\mathbb P}
\newcommand{\1}{\mathbf 1}
\newcommand{\Law}{\mathcal L}
\DeclareMathOperator{\rank}{rank}
\DeclareMathOperator{\diag}{diag}
\DeclareMathOperator{\im}{im}
\title{Polynomial Recovery in the Cubic Regular Stochastic Block Model}
\author{}
\date{}

\begin{document}
\maketitle

\begin{abstract}
We consider the regular stochastic block model with two planted classes
of size $N$, independent uniform simple $2$-regular graphs inside the
classes, and an independent uniform perfect matching between them.
There is a deterministic algorithm with running time bounded by a fixed
polynomial that recovers the planted partition, up to exchange of its
classes, with probability $1-O(N^{-2})$. The main estimate is the coarse
finite-field moment bound
$\E 2^{\dim_{\F_2}\ker(A-I)}=O(N^{24})$.
An exact configuration-model type formula and a short entropy
symmetrization prove this estimate. Binary linear algebra then reduces
recovery to a polynomial number of candidate checks with high probability.
We include a complete first-moment proof of the Boolean partition
rigidity needed to identify the planted candidate. The argument does
not require the real eigenspace at eigenvalue $1$ to be one-dimensional.
\end{abstract}

\section{Model and recovery theorem}

Let $V=V_+\sqcup V_-$, with $|V_+|=|V_-|=N$ and $N\geq3$.
Independently choose a uniform simple $2$-regular graph in each class
and a uniform perfect matching between the classes. Their union is a
simple cubic graph $G$, with adjacency matrix $A$ of order $n=2N$.
All asymptotic statements below are along the integers $N\geq3$; there
is no further parity restriction on $N$.
The planted sign vector $\sigma$ equals $1$ on $V_+$ and $-1$ on $V_-$.
It satisfies $A\sigma=\sigma$.

An \emph{equitable $(2,1)$ assignment} is a vector
$\tau\in\{1,-1\}^{n}$ for which every vertex has two neighbors with
its own sign and one with the opposite sign. On a cubic graph this is
equivalent to $A\tau=\tau$. Every such assignment is balanced:
symmetry and $A\1=3\1$ imply
\[
 3\langle\1,\tau\rangle
 =\langle\1,A\tau\rangle=\langle\1,\tau\rangle,
 \qquad\text{so}\qquad \langle\1,\tau\rangle=0.
\]
The assignments $\tau$ and $-\tau$ represent the same partition.

\begin{theorem}\label{thm:main}
In the model above there is a deterministic algorithm, given only $G$,
whose running time is bounded by a fixed polynomial in $n$ on every
input and which recovers the planted partition with probability
$1-O(N^{-2})$.
There is also a deterministic algorithm that lists all equitable
$(2,1)$ partitions on every input and has expected running time
bounded by a fixed polynomial under this model.
\end{theorem}

The proof separates two estimates. A finite-field first moment bounds
the number of candidates examined by the algorithm. A Boolean rigidity
estimate says that all passing candidates represent the planted
partition with high probability. The latter estimate is proved in
Appendix~\ref{app:rigidity}.

\section{An exact configuration-model moment}

First replace each internal graph by the degree-two configuration
model. Every vertex has two distinguished stubs, and a uniform perfect
pairing of the $2N$ stubs determines the internal multigraph.
Its adjacency matrix counts a loop twice on the diagonal. The two
internal pairings and the cross perfect matching are independent.

Relabel the second class using the cross matching. Relabeling
invariance and independence leave the two internal configuration
matrices $C_+,C_-$ independent with their original laws, and give
\[
 A=\begin{pmatrix}C_+&I\\I&C_-\end{pmatrix}.
\]
Conjugating $A-I$ by $\diag(I,-I)$ gives
\begin{equation}\label{eq:M}
 M=\begin{pmatrix}C_+-I&-I\\-I&C_--I\end{pmatrix}.
\end{equation}
This conjugation is invertible over every prime field. The relabeling
and conjugation are for analysis only; the algorithms will not need
the planted classes, matching, or signs.

Fix a prime $p$. A vector $(x,z)\in\F_p^N\times\F_p^N$ lies in
$\ker M$ precisely when
\begin{equation}\label{eq:local-equations}
 C_+x=x+z,\qquad C_-z=x+z.
\end{equation}
At coordinate $i$, let $a=x_i$ and $d=z_i$. Let $b,c$ be the $x$
labels at the partners of the two distinguished plus stubs at $i$,
and let $e,f$ be the $z$ labels at the partners of the two minus
stubs. The equations require
\[
 b+c=a+d=e+f.
\]
Thus an ordered local profile is determined by
$(a,d,b,e)\in\F_p^4$, with $c=a+d-b$ and $f=a+d-e$.
Let $n_{a,d,b,e}$ count the coordinates with each profile. Define
ordinary nonnegative integer tables
\begin{align*}
 m^+_{uv}
 &=\sum_{a,d,b,e}n_{a,d,b,e}\,
       \boldsymbol{1}_{a=u}
       \bigl(\boldsymbol{1}_{b=v}+\boldsymbol{1}_{a+d-b=v}\bigr),\\
 m^-_{uv}
 &=\sum_{a,d,b,e}n_{a,d,b,e}\,
       \boldsymbol{1}_{d=u}
       \bigl(\boldsymbol{1}_{e=v}+\boldsymbol{1}_{a+d-e=v}\bigr).
\end{align*}
The arithmetic in the subscripts is in $\F_p$; the indicator values
and table entries are integers. Call a table \emph{valid} if it is
symmetric and its diagonal entries are even. For such a table put
\[
 W(m)=\prod_{u<v}m_{uv}!\prod_u(m_{uu}-1)!!,
 \qquad (-1)!!=1,
\]
where any fixed ordering of $\F_p$ can be used in the first product.

\begin{proposition}\label{prop:exact}
In the degree-two configuration model,
\begin{equation}\label{eq:exact}
 \E|\ker_{\F_p}M|
 =\sum_{\substack{(n_{a,d,b,e}):\ \sum n_{a,d,b,e}=N\\
                         m^+,m^-\ \mathrm{valid}}}
   \frac{N!}{\prod_{a,d,b,e}n_{a,d,b,e}!}
   \frac{W(m^+)W(m^-)}{((2N-1)!!)^2}.
\end{equation}
\end{proposition}

\begin{proof}
Fix an assignment of profiles to the $N$ labeled coordinates. A plus
stub in class $(u,v)$ has host label $u$ and requests a partner whose
host label is $v$. For $u\ne v$, the stubs in class $(u,v)$ must be
matched bijectively to those in class $(v,u)$, giving $m_{uv}!$
choices if their sizes agree. The class $(u,u)$ must be paired
internally, giving $(m_{uu}-1)!!$ choices if its size is even. These
conditions are necessary and sufficient, so the compatible plus
pairings number $W(m^+)$. The independent minus count is $W(m^-)$.

There are $N!/\prod n_{a,d,b,e}!$ assignments of a fixed profile type.
Every pair of stub pairings and every kernel vector determine exactly
one ordered profile assignment. Conversely, a compatible pair of
pairings and a profile assignment give the kernel equations
\eqref{eq:local-equations}. Distinguished stubs ensure that no
additional factors of two are needed. Dividing the total count by
the number $((2N-1)!!)^2$ of pairing pairs proves the formula.
Loops and parallel edges are included by this count.
\end{proof}

\section{Entropy symmetrization and a polynomial moment}

Use natural logarithms and the convention $0\log0=0$. For a type in
\eqref{eq:exact}, put
\[
 T=\frac nN,\qquad Q^+=\frac{m^+}{2N},\qquad
 Q^-=\frac{m^-}{2N}.
\]
Regard $T$ as the law of random variables $(a,d,b,e)$, with
$c=a+d-b$ and $f=a+d-e$. Then
\[
 Q^+=\tfrac12\Law(a,b)+\tfrac12\Law(a,c),\qquad
 Q^-=\tfrac12\Law(d,e)+\tfrac12\Law(d,f).
\]
For a finite probability law $R$, denote its Shannon entropy by $H(R)$.

\begin{lemma}\label{lem:entropy}
Every law $T$ of $(a,d,b,e)\in\F_p^4$ satisfies
\begin{equation}\label{eq:entropy}
 \Phi(T):=H(T)-H(Q^+)-H(Q^-)\leq0.
\end{equation}
\end{lemma}

\begin{proof}
Consider the commuting involutions
\[
 s(a,d,b,e)=(a,d,a+d-b,e),\qquad
 t(a,d,b,e)=(a,d,b,a+d-e).
\]
Each is a bijection of the local profile space and leaves the two
averaged pair laws $Q^+,Q^-$ unchanged. Let $\overline T$ be the
average of $T,sT,tT,stT$. Entropy concavity and invariance under
bijections give $H(T)\leq H(\overline T)$. Under $\overline T$ the
two stubs in each class are exchangeable, so
\[
 Q^+=\Law_{\overline T}(a,b),\qquad
 Q^-=\Law_{\overline T}(d,e).
\]
Subadditivity for the two random pairs $(a,b)$ and $(d,e)$ now gives
\[
 \Phi(T)\leq H(\overline T)-H_{\overline T}(a,b)
                   -H_{\overline T}(d,e)
 =-I_{\overline T}\bigl((a,b);(d,e)\bigr)\leq0.
\]
The argument works in characteristic two and does not use the
pairing-validity conditions on $Q^+,Q^-$.
\end{proof}

\begin{proposition}\label{prop:moment-config}
For a fixed prime $p$, the configuration model satisfies
\begin{equation}\label{eq:moment-general}
 \E|\ker_{\F_p}M|
 \leq e^{p(p+1)+2}(2N+1)^{p(p+1)+2}(N+1)^{p^4}.
\end{equation}
In particular,
\begin{equation}\label{eq:moment-two-config}
 \E\,2^{\dim_{\F_2}\ker(A-I)}
 \leq e^8(2N+1)^8(N+1)^{16}=O(N^{24}).
\end{equation}
\end{proposition}

\begin{proof}
For every integer $r\geq0$, elementary factorial bounds give
\begin{equation}\label{eq:factorial}
 r\log r-r\leq\log(r!)
 \leq r\log r-r+1+\log(r+1).
\end{equation}
For positive $r$ these follow by comparing the sum of logarithms with
its integral; $r=0$ satisfies the displayed inequalities as well.
For even $r\geq0$, use
$(r-1)!!=r!/(2^{r/2}(r/2)!)$ to obtain
\[
 \log((r-1)!!)
 \leq\frac r2\log r-\frac r2+1+\log(r+1).
\]
Set $q=p(p+1)/2$. For a valid symmetric table $m$, with
$Q=m/(2N)$, summing these bounds gives
\begin{align*}
 \log W(m)
 &\leq\frac12\sum_{u,v}m_{uv}\log m_{uv}-N
                         +q\bigl(1+\log(2N+1)\bigr)\\
 &=N\log(2N)-N-NH(Q)+q\bigl(1+\log(2N+1)\bigr).
\end{align*}
Using \eqref{eq:factorial} in the opposite directions in
$(2N-1)!!=(2N)!/(2^NN!)$ gives
\[
 \log((2N-1)!!)
 \geq N\log(2N)-N-1-\log(N+1).
\]
Consequently
\begin{equation}\label{eq:pair-ratio}
 \frac{W(m)}{(2N-1)!!}
 \leq e^{q+1}(2N+1)^{q+1}e^{-NH(Q)}.
\end{equation}
The multinomial entropy bound is
$N!/\prod n!\leq e^{NH(T)}$. Combining it with
\eqref{eq:pair-ratio} and Lemma~\ref{lem:entropy}, every summand in
\eqref{eq:exact} is at most
$e^{2(q+1)}(2N+1)^{2(q+1)}$.
There are at most $(N+1)^{p^4}$ nonnegative profile-count tables.
This proves \eqref{eq:moment-general}. Taking $p=2$ and using the
invertible conjugation \eqref{eq:M} proves
\eqref{eq:moment-two-config}.
\end{proof}

\section{Conditioning to the uniform simple model}

\begin{lemma}\label{lem:simple}
A degree-two configuration on $N$ labeled vertices is simple with
probability tending to $e^{-3/4}$. Conditional on simplicity, its graph
is uniform among all simple labeled $2$-regular graphs.
\end{lemma}

\begin{proof}
Every simple $2$-regular graph has exactly $2^N$ distinguished-stub
realizations: at every vertex choose which of its two stubs is assigned
to each of its two incident edges. Thus conditioning gives the uniform
simple law.

For completeness, a simple $2$-regular graph is a collection of
undirected cycles of lengths at least three. The number of labeled
graphs divided by $N!$ is the coefficient of $z^N$ in
\[
 \exp\left(\sum_{\ell\geq3}\frac{z^\ell}{2\ell}\right)
 =(1-z)^{-1/2}e^{-z/2-z^2/4}.
\]
On the other hand,
\[
 \frac{(2N-1)!!}{2^NN!}
 =\frac{\binom{2N}{N}}{4^N}
 =[z^N](1-z)^{-1/2}\sim\frac1{\sqrt{\pi N}}.
\]
Multiplying $(1-z)^{-1/2}$ by the entire function
$g(z)=e^{-z/2-z^2/4}$ multiplies this coefficient asymptotic by
$g(1)=e^{-3/4}$. One can see this directly by writing the coefficient
as a convolution: for each fixed $j$, the ratio of the coefficients
at $N-j$ and $N$ tends to one; the coefficients of $g$ decay
exponentially, and the tail of the convolution is negligible.
Taking the ratio of the two counts proves the limit.
\end{proof}

\begin{corollary}\label{cor:simple-moment}
For the uniform simple model of Theorem~\ref{thm:main},
\begin{equation}\label{eq:simple-moment}
 \E\,2^{k}=O(N^{24}),\qquad
 k=\dim_{\F_2}\ker(A-I).
\end{equation}
In particular,
\begin{equation}\label{eq:cutoff-failure}
 \PP\{2^k>n^{27}\}=O(N^{-3}).
\end{equation}
\end{corollary}

\begin{proof}
The two internal configurations are independent, so their joint
simplicity probability tends to $e^{-3/2}>0$.
For the nonnegative random variable $2^k$, conditioning can increase
expectation by at most the reciprocal of this probability. Apply
Proposition~\ref{prop:moment-config} and then Markov's inequality.
\end{proof}

\begin{remark}
This conditioning step gives exactly the required uniform simple
$2$-regular law. No replacement by the adjacency matrix $P+P^{\mathsf T}$
of a single uniform permutation is used.
\end{remark}

\section{Recovery algorithms and complexity}

\subsection{Binary affine enumeration}

Write a candidate sign assignment as $\tau_i=(-1)^{z_i}$, with
$z\in\F_2^n$. The parity of the number of opposite-sign neighbors at
vertex $i$ is
\[
 \sum_{j\sim i}(z_i+z_j)=((A+I)z)_i\quad\text{in }\F_2,
\]
because the graph is cubic. Every equitable $(2,1)$ assignment therefore
satisfies
\begin{equation}\label{eq:affine}
 (A+I)z=\1\quad\text{over }\F_2.
\end{equation}
The parity equation permits either one or three opposite neighbors, so
its solutions must still be checked on the graph.

The bounded-time algorithm is as follows.
\begin{enumerate}
\item Solve \eqref{eq:affine} by binary Gaussian elimination. If it is
inconsistent, return failure. Otherwise write its solution set as
$z_0+\ker(A+I)$, and let $k=\dim_{\F_2}\ker(A+I)$.
\item If $2^k>n^{27}$, return failure.
\item Enumerate the affine solutions. For each solution, check that
every vertex has exactly one neighbor with a different label. Return
the first passing assignment, with a fixed global orientation.
If none passes, return failure.
\end{enumerate}

The planted assignment always satisfies \eqref{eq:affine}, so it is
among the enumerated candidates unless the cutoff triggers.
In fact the affine system is consistent for every simple cubic graph:
if $L=A+I$ and $u\in\ker_{\F_2}L$, then symmetry and the unit diagonal
give
\[
 0=u^{\mathsf T}Lu=\sum_i u_i=\1^{\mathsf T}u.
\]
Thus $\1\in(\ker L)^\perp=\im L$.

Elimination takes $O(n^3)$ binary operations using a dense matrix.
Use a binary Gray-code order to enumerate the coefficient vectors of
the affine space. Consecutive solutions differ by one kernel basis
vector, so updating the candidate takes $O(n)$ binary operations.
Checking it takes $O(n)$ operations on the cubic adjacency lists.
Hence the work is $O(n^3+n2^k)$ elementary operations. With the cutoff
this is $O(n^{28})$ in the usual RAM model, on every input. Accounting
for vertex indices of $O(\log n)$ bits introduces at most a logarithmic
factor in the corresponding indexed bit model. The memory use,
excluding streamed output, is $O(n^2)$ bits up to index storage.

Without the cutoff, the same enumeration always lists all equitable
assignments, and \eqref{eq:simple-moment} bounds its expected work by
\[
 O\bigl(n^3+n\E2^k\bigr)=O(n^{25})
\]
elementary operations. Keeping only assignments whose first vertex
has label zero removes the twofold duplication of partitions.

\begin{proof}[Proof of Theorem~\ref{thm:main}]
By Corollary~\ref{cor:simple-moment}, the cutoff triggers with
probability $O(N^{-3})$. By Proposition~\ref{prop:rigidity} in
Appendix~\ref{app:rigidity}, with probability $1-O(N^{-2})$ all passing
assignments are $\sigma$ and $-\sigma$. On the intersection of these
events the algorithm returns the planted partition. The complexity
claims follow from the bounds just established.
\end{proof}

\subsection{An exact spectral version}

There is also a direct rational eigenspace algorithm. Put
$r=\dim_{\mathbb Q}\ker(A-I)$. Reduction of integer minors modulo two
gives
\[
 \rank_{\F_2}(A-I)\leq\rank_{\mathbb Q}(A-I),\qquad r\leq k.
\]
Thus $\E2^r=O(N^{24})$ as well. Compute a rational basis matrix
$U\in\mathbb Q^{n\times r}$ and choose a set $J$ of $r$ rows for which
$U_J$ is invertible. For every $\varepsilon\in\{1,-1\}^r$, form
\[
 v=UU_J^{-1}\varepsilon.
\]
Keep exactly those $v$ all of whose coordinates belong to $\{1,-1\}$.
This lists precisely all Boolean eigenvectors at eigenvalue $1$:
every vector in the eigenspace is uniquely determined by its
coordinates on $J$. Apply the cutoff $2^r\leq n^{27}$ to obtain a
bounded-time recovery algorithm with the same probability estimate.

This is a polynomial algorithm in the bit model, not just in a unit-cost
rational arithmetic model. Every row of $A-I$ has four entries of
magnitude one, so every minor has absolute value at most $2^n$ by
Hadamard's inequality. Fraction-free elimination therefore constructs
a basis with polynomial bit length. More explicitly, Cramer's rule
gives a basis with a common denominator and numerators having $O(n)$
bits. Writing that basis as $U=V/\Delta$, the reconstruction matrix is
$VV_J^{-1}$. Its entries have numerator and denominator bit lengths
$O(n^2+n\log n)$, again by determinant bounds. Basis construction,
each reconstruction, and each sign check consequently take polynomial
bit complexity. The cutoff bounds the number of reconstructions by a
fixed polynomial; without it, their expected number is polynomial.

The spectral version needs only logarithmic real multiplicity with
high probability. Neither version asserts that
$\dim_{\mathbb R}\ker(A-I)=1$ with high probability, and the polynomial
exponents here are coarse existence bounds rather than practical
running-time estimates.

\appendix
\section{A complete Boolean rigidity bound}\label{app:rigidity}

\begin{proposition}\label{prop:rigidity}
In the uniform simple model, the probability that there is an equitable
$(2,1)$ partition different from the planted partition is $O(N^{-2})$.
Equivalently, with probability $1-O(N^{-2})$, the only vectors in
$\{1,-1\}^{2N}$ satisfying $A\tau=\tau$ are $\sigma$ and $-\sigma$.
\end{proposition}

\subsection{Marked vertices and an exact first moment}

Given an alternative assignment, mark precisely the vertices where its
sign differs from $\sigma$. Let $S_+\subseteq V_+$ and
$S_-\subseteq V_-$ be the marked sets. At a vertex let $r$ be the
number of its two internal edges whose endpoints have different marks,
and let $m\in\{0,1\}$ indicate whether its original cross-matching
edge has differently marked endpoints. In the alternative assignment,
the number of opposite-sign neighbors is
\[
 r+1-m.
\]
Consequently the alternative is equitable precisely when
\begin{equation}\label{eq:r-m}
 r=m\quad\text{at every vertex}.
\end{equation}
In particular, the internal edges crossing a marked/unmarked cut form
a matching. A vertex belongs to this internal boundary matching exactly
when its cross-matching edge joins different marks.

Balance of the alternative assignment implies
$|S_+|=|S_-|=k$. Replacing the alternative by its negative replaces $k$
by $N-k$, so any alternative partition has a representative with
$1\leq k\leq N/2$. In a simple internal graph $k=1$ is impossible:
the sole marked vertex has two distinct unmarked internal neighbors,
which contradicts $r=m\leq1$.

Work temporarily in the unconditioned degree-two configuration model,
retaining an independent uniform cross perfect matching. Fix marked
sets of size $k$ in both classes. In one internal block, the probability
that the boundary is a matching with exactly $t$ edges is
\begin{equation}\label{eq:Pt}
 P_t=\binom kt\binom{N-k}t 2^{2t}t!
 \frac{(2k-t-1)!!\,(2(N-k)-t-1)!!}{(2N-1)!!},
\end{equation}
where $t$ is even and $0\leq t\leq k$.
Indeed, choose $t$ boundary vertices on each side, select one of their
two stubs at each vertex, pair those stubs across the cut, and pair
the remaining stubs within their sides. This gives exactly the
numerator in \eqref{eq:Pt}.

For \eqref{eq:r-m} to hold, the two internal boundary sizes must agree.
Conditional on their common size $t$, the probability that the cross
matching is compatible is
\begin{equation}\label{eq:Qt}
 Q_t=\frac{(t!)^2(k-t)!(N-k-t)!}{N!}.
\end{equation}
The four types in each block are marked boundary, unmarked boundary,
marked interior, and unmarked interior. The cross matching must
interchange the two boundary types and preserve the two interior types,
which proves \eqref{eq:Qt}.

Thus the expected number of alternatives with overlap size $k$ is
\begin{equation}\label{eq:Ek}
 E_k=\binom Nk^2\sum_{\substack{0\leq t\leq k\\t\ \mathrm{even}}}
                         P_t^2Q_t.
\end{equation}
The formula is an exact first-moment count before conditioning on
simplicity. We will prove
\begin{equation}\label{eq:rigidity-sum}
 \sum_{k=2}^{\lfloor N/2\rfloor}E_k=O(N^{-2}).
\end{equation}

\subsection{Small overlaps}

For $2\leq k\leq N/8$, cancel
$(2(N-k)-t-1)!!$ against the denominator in \eqref{eq:Pt}.
There are $k+t/2$ removed denominator factors, each at least $N$.
Using
\[
 \binom kt\leq2^k,\qquad
 \binom{N-k}t t!\leq N^t,\qquad
 (2k-t-1)!!\leq(2k)^{k-t/2}
\]
gives
\begin{equation}\label{eq:Pt-small}
 P_t\leq16^k(k/N)^{k-t/2}.
\end{equation}
Also $N-k-t\geq N/2$, so
\begin{equation}\label{eq:Qt-small}
 Q_t\leq4^k k^{k+t}N^{-k-t}.
\end{equation}
For example, bound the numerator factorials in \eqref{eq:Qt} by
$k^{2t}k^{k-t}$ and the quotient $N!/(N-k-t)!$ below by
$(N/2)^{k+t}$. Combining \eqref{eq:Pt-small},
\eqref{eq:Qt-small}, and $\binom Nk\leq(eN/k)^k$ gives, for each
summand $E_{k,t}=\binom Nk^2P_t^2Q_t$,
\[
 E_{k,t}\leq(C_0k/N)^k,\qquad C_0=1024e^2.
\]
Set $\delta=1/(100C_0)$. There are at most $k+1$ values of $t$, and
the ratio between consecutive quantities
$(k+1)(C_0k/N)^k$ is at most
\[
 \frac32 e C_0\frac{k+1}{N}.
\]
For all sufficiently large $N$ this is less than $1/2$ throughout
$2\leq k\leq\delta N$. The first term is $O(N^{-2})$, so
\begin{equation}\label{eq:small-overlaps}
 \sum_{k=2}^{\lfloor\delta N\rfloor}E_k=O(N^{-2}).
\end{equation}

\subsection{Macroscopic overlaps}

Write $\alpha=k/N$, $\beta=t/N$, and $s(x)=x\log x$, with $s(0)=0$.
Substitution of the factorial formulas into $E_{k,t}$, using
\eqref{eq:factorial} and the double-factorial identity, gives uniformly
for $\delta\leq\alpha\leq1/2$ and $0\leq\beta\leq\alpha$,
\begin{equation}\label{eq:F-approx}
 \log E_{k,t}=NF(\alpha,\beta)+O(\log(N+1)),
\end{equation}
where
\begin{align}
 F(\alpha,\beta)
 ={}&-s(\alpha-\beta)-s(1-\alpha-\beta)
      +s(2\alpha-\beta)+s(2-2\alpha-\beta)
      +(4\beta-2)\log2.\label{eq:F}
\end{align}
The error is uniform even when factorial arguments vanish, since
\eqref{eq:factorial} is uniform down to zero and only a fixed number of
factorials occur.

We prove that the maximum of $F$ on this range is strictly negative.
For fixed $\alpha\in(0,1/2]$, in the interior of the $\beta$ interval,
\[
 \frac{\partial F}{\partial\beta}
 =\log\frac{16(\alpha-\beta)(1-\alpha-\beta)}
                {(2\alpha-\beta)(2-2\alpha-\beta)}.
\]
Its derivative is strictly negative, because
\[
 -\frac1{\alpha-\beta}-\frac1{1-\alpha-\beta}
 +\frac1{2\alpha-\beta}+\frac1{2-2\alpha-\beta}<0.
\]
The first derivative is $\log4>0$ at $\beta=0$ and tends to
$-\infty$ at the other endpoint. The unique maximizer is therefore
\[
 b(\alpha)=\frac{7-\sqrt{49-180\alpha(1-\alpha)}}{15},
 \qquad 0<b(\alpha)\leq\frac13,
\]
with equality in the upper bound only at $\alpha=1/2$.
Let $f(\alpha)=F(\alpha,b(\alpha))$. Put
\[
 u=\alpha-b,\quad v=1-\alpha-b,\quad
 x=2u+b,\quad y=2v+b.
\]
The stationarity equation is $xy=16uv$, and the envelope derivative is
$f'(\alpha)=\log(x^2v/(y^2u))$.
The difference determining its sign factors as
\[
 x^2v-y^2u=(u-v)(4uv-b^2),\qquad
 4uv-b^2=2b(1/3-b),
\]
where the second identity follows by expanding $xy=16uv$ and using
$u+v=1-2b$. For $0<\alpha<1/2$, one has $u-v=2\alpha-1<0$ and
$0<b<1/3$. Thus $f'(\alpha)<0$. Finally $f(\alpha)\to0$ as
$\alpha\downarrow0$, by \eqref{eq:F}. Consequently
\[
 F(\alpha,\beta)\leq f(\alpha)\leq f(\delta)<0
 \quad\text{for}\quad
 \delta\leq\alpha\leq1/2,\quad0\leq\beta\leq\alpha.
\]
Equation \eqref{eq:F-approx}, followed by a sum over at most $N^2$
pairs $(k,t)$, proves an exponentially small upper bound for all
macroscopic overlaps. Together with \eqref{eq:small-overlaps}, this
proves \eqref{eq:rigidity-sum}.

Finally, $k=1$ is impossible in the simple model, and conditioning both
internal configurations to be simple costs a bounded factor by
Lemma~\ref{lem:simple}. Thus \eqref{eq:rigidity-sum} also bounds the
probability of an alternative partition in the simple model, proving
Proposition~\ref{prop:rigidity}.

\end{document}
